% Luc Destombe --- licence CC BY-NC-SA 4.0
% https://creativecommons.org/licenses/by-nc-sa/4.0/deed.fr
% Fichier autonome : compiler avec pdflatex, deux fois (signets, nombre de pages).
\documentclass[11pt,a4paper]{article}
\makeatletter
\newif\iflycee@unecolonne
\lycee@unecolonnetrue
\newif\iflycee@palatino
\lycee@palatinotrue

\RequirePackage[utf8]{inputenc}
\RequirePackage[T1]{fontenc}
\RequirePackage[french]{babel}
\RequirePackage{etoolbox}

\newcommand{\ShorthandsOff}{\@ifundefined{shorthandoff}{}{\shorthandoff{;:!?}}}
\newcommand{\ShorthandsOn}{\@ifundefined{shorthandon}{}{\shorthandon{;:!?}}}
\ShorthandsOff
\AtBeginDocument{\ShorthandsOn}
\AtBeginEnvironment{tikzpicture}{\ShorthandsOff}

\RequirePackage{geometry}
\RequirePackage{fancyhdr}
\RequirePackage{lastpage}
\RequirePackage{multicol}
\RequirePackage{array}
\RequirePackage{enumitem}
\RequirePackage{xcolor}
\RequirePackage{needspace}

\RequirePackage{amsmath,amssymb}
\RequirePackage{mathpazo}
\DeclareMathSizes{10.95}{10.95}{8}{5.5}
\begingroup\catcode`\_=\active \catcode`\^=\active
\gdef_#1{\sb{\scriptscriptstyle #1}}
\gdef^#1{\sp{\scriptscriptstyle\lycee@moinsepais #1}}
\endgroup
\newcommand\lycee@moins{\mathbin{%
  \setbox\z@\hbox{$\m@th\scriptscriptstyle\lycee@moinsorig$}%
  \dimen@=.55\wd\z@ \advance\dimen@-.5pt
  \hbox to.75\wd\z@{\hss$\m@th\scriptscriptstyle\vcenter{\hbox{%
    \tikz\draw[line width=.5pt,line cap=round](0,0)--(\the\dimen@,0);}}$\hss}}}
\begingroup\lccode`\~=`\- \lowercase{\endgroup\let~}\lycee@moins
\newcommand\lycee@moinsepais{\mathcode`\-="8000 }
\AtBeginDocument{\mathchardef\lycee@moinsorig=\mathcode`\-\relax}
\AtBeginDocument{\catcode`\_=12 \mathcode`\_="8000
  \catcode`\^=12 \mathcode`\^="8000 }
\newcommand\lycee@exposants{%
  \fontdimen13\textfont2=.48\fontdimen6\textfont2
  \fontdimen14\textfont2=.48\fontdimen6\textfont2
  \fontdimen15\textfont2=.48\fontdimen6\textfont2
  \fontdimen13\scriptfont2=.48\fontdimen6\scriptfont2
  \fontdimen14\scriptfont2=.48\fontdimen6\scriptfont2
  \fontdimen15\scriptfont2=.48\fontdimen6\scriptfont2}
\every@math@size\expandafter{\the\every@math@size\lycee@exposants}
\newlength{\retraitcalculs}
\setlength{\retraitcalculs}{2em}

\RequirePackage{tikz}
\usetikzlibrary{arrows.meta,calc,babel,shapes.misc}
\RequirePackage[most]{tcolorbox}
\tcbuselibrary{skins,breakable}

\definecolor{coulDef}{HTML}{1F4E79}
\definecolor{coulProp}{HTML}{7030A0}
\definecolor{coulRem}{HTML}{C55A11}
\definecolor{coulEx}{HTML}{2E7D32}
\definecolor{coulExo}{HTML}{B03A2E}
\definecolor{coulPart}{HTML}{1F4E79}
\definecolor{coulMeth}{HTML}{1B6E4F}
\definecolor{coulCourbe}{HTML}{0B5ED7}
\definecolor{coulCourbeB}{HTML}{C00000}
\definecolor{coulCourbeC}{HTML}{2E7D32}
\definecolor{coulGrille}{HTML}{8C8C8C}
\definecolor{coulRep}{HTML}{1B6E4F}
\definecolor{coulBar}{HTML}{2E7D32}

\newcommand{\R}{\mathbb{R}}
\newcommand{\Rs}{\mathbb{R}^{*}}
\renewcommand{\le}{\leqslant}
\renewcommand{\ge}{\geqslant}
\renewcommand{\approx}{\simeq}
\newcommand{\vect}[1]{\overrightarrow{#1}}
\newcommand{\norme}[1]{\left\lVert #1 \right\rVert}
\newcommand{\intervalle}[2]{\left[#1\,;#2\right]}
\RequirePackage{cancel}
\renewcommand{\CancelColor}{\color{coulExo}}
\newcommand{\fois}[1]{\times{\color{coulExo}#1}}
\newcommand{\barre}[1]{\cancel{#1}}

\tikzset{grille/.style={coulGrille,line width=0.4pt}}
\newcommand{\quadrillage}[4]{\draw[grille,step=1] (#1,#2) grid (#3,#4);}
\tikzset{
  croix/.style={cross out,draw,line width=0.6pt,inner sep=0pt,minimum size=4pt},
  croixdroite/.style={croix,rotate=45,line width=0.8pt},
}
\newcommand{\pt}[3]{\path (#1) node[croix]{} node[#2,font=\small,inner sep=2.5pt]{$#3$};}

\newcommand{\lycee@repere}[4]{%
  \draw[grille,step=1] (#1,#2) grid (#3,#4);
  \draw[-{Stealth[length=2mm]},thick] (#1,0)--({#3+0.4},0) node[below left=-1pt]{$x$};
  \draw[-{Stealth[length=2mm]},thick] (0,#2)--(0,{#4+0.4}) node[below left=-1pt]{$y$};
  \node[below left,font=\scriptsize,inner sep=1.5pt] at (0,0) {$0$};}
\newcommand{\lycee@gradx}[1]{\foreach \k/\l in {#1}{%
  \draw (\k,0.07)--(\k,-0.07) node[below,font=\scriptsize,inner sep=1.5pt]{$\l$};}}
\newcommand{\lycee@grady}[1]{\foreach \k/\l in {#1}{%
  \draw (0.07,\k)--(-0.07,\k) node[left,font=\scriptsize,inner sep=1.5pt]{$\l$};}}
\newcommand{\lycee@intff}[2]{\left[#1\,;#2\right]}
\newcommand{\lycee@intfo}[2]{\left[#1\,;#2\right[}
\newcommand{\lycee@intof}[2]{\left]#1\,;#2\right]}
\newcommand{\lycee@intoo}[2]{\left]#1\,;#2\right[}
\newcommand{\lycee@ens}[1]{\left\{#1\right\}}
\AtBeginDocument{%
  \providecommand{\repere}{\lycee@repere}%
  \providecommand{\gradx}{\lycee@gradx}%
  \providecommand{\grady}{\lycee@grady}%
  \providecommand{\Cf}{\mathcal{C}\sb{\scriptscriptstyle f}}%
  \providecommand{\Cg}{\mathcal{C}\sb{\scriptscriptstyle g}}%
  \providecommand{\intff}{\lycee@intff}%
  \providecommand{\intfo}{\lycee@intfo}%
  \providecommand{\intof}{\lycee@intof}%
  \providecommand{\intoo}{\lycee@intoo}%
  \providecommand{\ens}{\lycee@ens}}

\newlist{questions}{enumerate}{3}
\setlist[questions,1]{label=\arabic*\textdegree),leftmargin=*,itemsep=2pt,topsep=3pt}
\setlist[questions,2]{label=\alph*),leftmargin=*,itemsep=1pt,topsep=2pt}
\setlist[questions,3]{label=\roman*),leftmargin=*,itemsep=1pt,topsep=1pt}
\setlist[itemize]{label=\textbullet,leftmargin=*,itemsep=2pt,topsep=3pt}

\newcommand{\besoinplace}[1]{\par\penalty\@M\begingroup
  \dimen@=#1\relax\dimen@ii=\pagegoal\advance\dimen@ii-\pagetotal
  \ifdim\dimen@>\dimen@ii\ifdim\dimen@ii>\z@\vfil\fi\break\fi\endgroup}

\newcommand{\lycee@pied}{\thepage/\pageref{LastPage}}
\newcommand{\entete}[2]{%
  \pagestyle{fancy}\fancyhf{}%
  \fancyhead[L]{\footnotesize\color{coulPart}\textbf{#1}}%
  \fancyhead[R]{\footnotesize\color{coulPart}\textbf{#2}}%
  \fancyfoot[C]{\footnotesize\lycee@pied}%
  \renewcommand{\headrulewidth}{0.4pt}%
  \renewcommand{\headrule}{\hbox to\headwidth{%
    \color{coulPart!40}\leaders\hrule height \headrulewidth\hfill}}}

\RequirePackage{ccicons}
\newcommand{\lycee@licence}{{\scriptsize\color{gray}%
  \copyright\ Luc Destombe --- \ccbyncsa\ CC BY-NC-SA 4.0}}
\newif\iflycee@licence \lycee@licencetrue
\newcommand{\sanslicence}{\lycee@licencefalse}
\AtBeginDocument{\iflycee@licence\AddToHookNext{shipout/foreground}{%
  \put(0,0){\rlap{\kern\dimexpr1in+\hoffset+\oddsidemargin+\textwidth\relax
    \llap{\raisebox{-\dimexpr1in+\voffset+\topmargin+\headheight+\headsep
      +\textheight+\footskip\relax}{\lycee@licence}}}}}\fi}

\newcommand{\formulesserrees}{%
  \setlength{\abovedisplayskip}{3pt}\setlength{\belowdisplayskip}{3pt}%
  \setlength{\abovedisplayshortskip}{2pt}\setlength{\belowdisplayshortskip}{3pt}}

\setlength{\parindent}{0pt}

\RequirePackage[implicit=false,hidelinks,pdfencoding=auto,psdextra,bookmarksopen,
  bookmarksopenlevel=1,pdfcreator={},pdfproducer={}]{hyperref}
\RequirePackage{bookmark}
\hypersetup{pdfauthor={Luc Destombe},pdfkeywords={CC BY-NC-SA 4.0}}
\newcounter{lycee@signet}
\newcommand{\signet}[3][]{\stepcounter{lycee@signet}%
  \edef\lycee@dest{\if\relax\detokenize{#1}\relax
    signet.\arabic{lycee@signet}\else#1\fi}%
  \hypertarget{\lycee@dest}{}\bookmark[level=#2,dest=\lycee@dest]{#3}}
\pdfstringdefDisableCommands{%
  \def\R{R}\def\vect#1{#1}\def\Cf{Cf}\def\Cg{Cg}\def\fois#1{×#1}%
  \def\barre#1{#1}\def\intervalle#1#2{[#1;#2]}\def\intff#1#2{[#1;#2]}%
  \def\textsuperscript#1{#1}\def\dfrac#1#2{#1/#2}\def\frac#1#2{#1/#2}%
  \def\vec#1{#1⃗}}%
\RequirePackage[official]{eurosym}

\tcbset{exercice corrige/.style={
  enhanced,breakable,before upper=\raggedright,
  colback=white,colframe=coulExo,
  boxrule=0.8pt,arc=2pt,
  left=6pt,right=6pt,top=8pt,bottom=5pt,
  before skip=9pt,after skip=4pt,
  coltitle=white,fonttitle=\bfseries\small,
  attach boxed title to top left={xshift=8pt,yshift=-\tcboxedtitleheight/2},
  boxed title style={colback=coulExo,arc=2pt,boxrule=0pt}}}
\newcounter{exo}
\newtcolorbox{exercice}[1][]{exercice corrige,
  phantom={\signet[exercice-\arabic{exo}]{2}{Exercice \arabic{exo}}},
  title={Exercice \theexo\ifx\relax#1\relax\else{} --- #1\fi}}
\newenvironment{exo}[1][\relax]{\refstepcounter{exo}\begin{exercice}[#1]}{\end{exercice}}

\geometry{top=1.8cm,bottom=1cm,left=1cm,right=1cm,headheight=14pt,footskip=0.6cm}
\renewcommand{\lycee@pied}{\thepage}

\renewcommand{\theexo}{\ifnum\value{exo}<10 0\fi\arabic{exo}}
\tcbset{exercice corrige/.style={
  enhanced,breakable,before upper=\raggedright,
  colback=white,colframe=coulExo,
  boxrule=0.9pt,arc=2pt,
  left=8pt,right=8pt,top=8pt,bottom=6pt,
  before skip=8pt,after skip=6pt,
  coltitle=white,fonttitle=\bfseries,
  attach boxed title to top left={xshift=8pt,yshift=-\tcboxedtitleheight/2},
  boxed title style={colback=coulExo,arc=2pt,boxrule=0pt}}}

\newcommand{\entetefiche}[2]{%
  \begin{center}
    \begin{tcolorbox}[enhanced,width=0.92\linewidth,colback=white,colframe=coulPart,
        boxrule=1.4pt,arc=3pt,halign=center,top=6pt,bottom=6pt]
      {\large\bfseries\color{coulPart} CORRIGÉ --- FICHE D'EXERCICES}\\[3pt]
      {\normalsize\bfseries #1}\\[2pt]
      {\normalsize #2}
    \end{tcolorbox}
  \end{center}}

\newcounter{partie}
\newcommand{\partie}[1]{%
  \stepcounter{partie}%
  \besoinplace{8\baselineskip}\vspace{4pt}%
  \begin{tcolorbox}[enhanced,colback=coulPart!8,colframe=coulPart,boxrule=0pt,
      leftrule=3pt,arc=0pt,left=6pt,right=4pt,top=3pt,bottom=3pt,
      before skip=6pt,after skip=2pt]
    \bfseries\color{coulPart}\Roman{partie} --- #1\signet{1}{\Roman{partie} --- #1}
  \end{tcolorbox}}
\newcommand{\partiesansnum}[1]{%
  \besoinplace{8\baselineskip}\vspace{4pt}%
  \begin{tcolorbox}[enhanced,colback=coulPart!8,colframe=coulPart,boxrule=0pt,
      leftrule=3pt,arc=0pt,left=6pt,right=4pt,top=3pt,bottom=3pt,
      before skip=6pt,after skip=2pt]
    \bfseries\color{coulPart}#1\signet{1}{#1}
  \end{tcolorbox}}

\setlist[questions,1]{label=\arabic*\textdegree),leftmargin=*,itemsep=3pt,topsep=3pt}
\setlist[questions,2]{label=\alph*),leftmargin=*,itemsep=2pt,topsep=2pt}
\setlist[questions,3]{label=\roman*),leftmargin=*,itemsep=1pt,topsep=1pt}
\newlist{expr}{enumerate}{1}
\setlist[expr]{label={\Alph*\ $=$},leftmargin=*,itemsep=2pt,topsep=3pt}

\newcommand{\res}[1]{{\color{coulRep}#1}}
\newcommand{\rem}[1]{\par\smallskip{\small\itshape\color{black!60}$\rightarrow$ #1}}
\newenvironment{calculs}{\par\smallskip\noindent$\displaystyle\begin{aligned}[t]}%
  {\end{aligned}$\par\smallskip}
\newcommand{\justif}[1]{\quad\text{\small\itshape\color{black!60}#1}}
\newcommand{\cas}[1]{\par\smallskip\noindent\textbf{\color{coulExo}#1}\nobreak}
\newcommand{\explique}[2]{\par\smallskip\noindent
  \begin{minipage}[t]{0.6\linewidth}\raggedright #1\end{minipage}\hfill
  \begin{minipage}[t]{0.37\linewidth}\raggedright\leavevmode
    {\small\itshape\color{black!60}#2}\end{minipage}\par}

\setlength{\parskip}{2pt}

\makeatother
\entete{Chapitre 7 --- Pourcentages --- corrigé}{Seconde}
\setlength{\parindent}{0pt}
\setlength{\jot}{5pt}
\let\casligne\cas
\renewcommand{\cas}[1]{\casligne{#1}\par\nopagebreak}
\renewcommand{\res}[1]{\leavevmode{\color{coulRep}#1}}

\begin{document}

\entetefiche{Chapitre 7 --- Pourcentages}{Seconde}

\partie{Appliquer et calculer un pourcentage}

\begin{exo}[Prendre un pourcentage]
\explique{Prendre $t\,\%$ d'une quantité, c'est la multiplier par
$\dfrac{t}{100}$.}{définition 1 du cours}
\cas{a)}
\begin{calculs}
10\,\%\text{ de }340 &= 0{,}1\times 340 = \res{34} &&\justif{$10\,\%$ : le dixième}\\
25\,\%\text{ de }48 &= 0{,}25\times 48 = \res{12} &&\justif{$25\,\%$ : le quart}\\
50\,\%\text{ de }3\,600 &= 0{,}5\times 3\,600 = \res{1\,800} &&\justif{$50\,\%$ : la moitié}
\end{calculs}
\cas{b)}
\begin{calculs}
15\,\%\text{ de }260 &= 0{,}15\times 260 = \res{39} &&\justif{$\frac{15}{100}=0{,}15$}\\
8\,\%\text{ de }75 &= 0{,}08\times 75 = \res{6} &&\justif{$\frac{8}{100}=0{,}08$, pas $0{,}8$}\\
3{,}5\,\%\text{ de }400 &= 0{,}035\times 400 = \res{14} &&\justif{$\frac{3{,}5}{100}=0{,}035$}
\end{calculs}
\cas{c)}
\begin{calculs}
200\,\%\text{ de }17 &= 2\times 17 = \res{34} &&\justif{$200\,\%$ : le double}\\
120\,\%\text{ de }45 &= 1{,}2\times 45 = \res{54} &&\justif{plus de $100\,\%$ : plus que $45$}\\
0{,}5\,\%\text{ de }1\,800 &= 0{,}005\times 1\,800 = \res{9} &&\justif{$\frac{0{,}5}{100}=0{,}005$}
\end{calculs}
\rem{Erreur fréquente : écrire $8\,\%=0{,}8$. On divise par $100$ : la virgule recule de
deux rangs, $8\,\%=0{,}08$.}
\end{exo}

\begin{exo}[Appliquer un pourcentage]
\cas{a)}
\begin{calculs}
42\,\%\text{ de }1\,150 &= 0{,}42\times 1\,150 &&\justif{on prend $42\,\%$ de l'effectif}\\
42\,\%\text{ de }1\,150 &= 483
\end{calculs}
\explique{Le lycée compte \res{$483$ demi-pensionnaires}.}{on conclut par une phrase}
\cas{b)}
\begin{calculs}
\text{réduction} &= 0{,}15\times 420 &&\justif{la réduction est $15\,\%$ du prix}\\
\text{réduction} &= 63\\
\text{prix payé} &= 420-63 &&\justif{on retire la réduction}\\
\text{prix payé} &= 357
\end{calculs}
\explique{La réduction est de \res{$63$~€} ; le téléphone est payé
\res{$357$~€}.}{une phrase par question}
\cas{c)}
\explique{$65\,\%$ du fichier est téléchargé : il reste $100\,\%-65\,\%=35\,\%$ du
fichier.}{on cherche ce qui \emph{reste}}
\begin{calculs}
35\,\%\text{ de }360 &= 0{,}35\times 360\\
35\,\%\text{ de }360 &= 126
\end{calculs}
\explique{Il reste \res{$126$~Mo} à télécharger.}{autre voie : $0{,}65\times 360=234$~Mo
téléchargés, puis $360-234=126$}
\end{exo}

\begin{exo}[Calculer un pourcentage]
\explique{On divise l'effectif concerné par l'effectif total, puis on écrit le quotient en
centièmes.}{méthode 2 du cours}
\cas{a)}
\explique{La classe compte $18+7=25$ élèves qui ont répondu.}{l'effectif total n'est pas
écrit : on le calcule}
\begin{calculs}
\dfrac{18}{25} &= 0{,}72 &&\justif{les bonnes réponses, divisées par le total}\\
\dfrac{18}{25} &= 72\,\% &&\justif{$0{,}72=\frac{72}{100}$}
\end{calculs}
\explique{\res{$72\,\%$} des élèves ont répondu juste.}{diviser par $7$ serait une erreur :
$7$ n'est pas le total}
\cas{b)}
\begin{calculs}
\dfrac{40}{320} &= 0{,}125 &&\justif{les embauches, divisées par l'effectif de départ}\\
\dfrac{40}{320} &= 12{,}5\,\%
\end{calculs}
\explique{Les embauches représentent \res{$12{,}5\,\%$} de l'effectif de départ.}{on
conclut par une phrase}
\cas{c)}
\begin{calculs}
\dfrac{288}{1\,600} &= 0{,}18 &&\justif{les courses, divisées par le salaire}\\
\dfrac{288}{1\,600} &= 18\,\%
\end{calculs}
\explique{Les courses représentent \res{$18\,\%$} du salaire.}{on conclut par une phrase}
\end{exo}

\begin{exo}[Comparer des pourcentages]
\cas{a)}
\begin{calculs}
\dfrac{12}{30} &= 0{,}4 &&\justif{classe A}\\
\dfrac{12}{30} &= 40\,\%\\
\dfrac{13}{35} &\approx 0{,}3714 &&\justif{classe B (calculatrice : $0{,}371428\ldots$)}\\
\dfrac{13}{35} &\approx 37{,}1\,\% &&\justif{arrondi au dixième}\\
\dfrac{11}{28} &\approx 0{,}3929 &&\justif{classe C (calculatrice : $0{,}392857\ldots$)}\\
\dfrac{11}{28} &\approx 39{,}3\,\%
\end{calculs}
\explique{A : \res{$40\,\%$} ; B : environ \res{$37{,}1\,\%$} ; C : environ
\res{$39{,}3\,\%$}.}{}
\cas{b)}
\explique{Rangement : \res{A, puis C, puis B}. La classe B a le plus d'élèves de théâtre
($13$), mais elle est la dernière : \res{non}.}{on compare les pourcentages, pas les
effectifs : les classes n'ont pas la même taille}
\end{exo}

\begin{exo}[Lire un tableau]
\cas{a)}
\begin{calculs}
\dfrac{504}{1\,200} &= 0{,}42 = \res{42\,\%} &&\justif{bus}\\
\dfrac{156}{1\,200} &= 0{,}13 = \res{13\,\%} &&\justif{vélo}\\
\dfrac{300}{1\,200} &= 0{,}25 = \res{25\,\%} &&\justif{à pied}\\
\dfrac{240}{1\,200} &= 0{,}2 = \res{20\,\%} &&\justif{voiture}
\end{calculs}
\explique{Vérification : $42+13+25+20=100$.}{les modes ne se recoupent pas : les
pourcentages font $100\,\%$}
\cas{b)}
\explique{\res{$\dfrac{300}{1\,200}$ est la proportion des élèves du lycée qui viennent à
pied : un quart d'entre eux, soit $25\,\%$.}}{population de référence : les $1\,200$
élèves ; sous-population : ceux qui viennent à pied}
\end{exo}

\begin{exo}[La TVA]
\cas{a)}
\begin{calculs}
\text{TVA} &= 0{,}2\times 240 &&\justif{$20\,\%$ du prix HT}\\
\text{TVA} &= 48\\
\text{prix TTC} &= 240+48 &&\justif{on ajoute la TVA au prix HT}\\
\text{prix TTC} &= 288
\end{calculs}
\explique{La TVA est de \res{$48$~€} ; le bureau coûte \res{$288$~€} TTC.}{une phrase par
question}
\cas{b)}
\begin{calculs}
\text{TVA} &= 0{,}2\times 85\\
\text{TVA} &= 17\\
\text{prix TTC} &= 85+17\\
\text{prix TTC} &= 102
\end{calculs}
\explique{La TVA est de \res{$17$~€} ; la chaise coûte \res{$102$~€} TTC.}{dans la partie IV,
on le fera en une multiplication : $85\times 1{,}2=102$}
\end{exo}

\partie{Proportions}

\begin{exo}[Effectif, proportion, effectif total]
\cas{a)}
\explique{Population de référence : les $320$ licenciés ; sous-population : les
mineurs.}{on repère d'abord les deux populations}
\begin{calculs}
n &= p\times N &&\justif{on cherche l'effectif de la sous-population}\\
n &= 0{,}35\times 320\\
n &= 112
\end{calculs}
\explique{Le club compte \res{$112$ licenciés mineurs}.}{on conclut par une phrase}
\cas{b)}
\explique{Population de référence : les $320$ licenciés ; sous-population : les $144$
femmes.}{}
\begin{calculs}
p &= \dfrac{n}{N} &&\justif{on cherche une proportion}\\
p &= \dfrac{144}{320}\\
p &= 0{,}45\\
p &= 45\,\%
\end{calculs}
\explique{Les femmes représentent \res{$45\,\%$} des licenciés.}{on conclut par une phrase}
\cas{c)}
\explique{Population de référence : les $144$ femmes ; sous-population : les $72$
mineures.}{la population de référence a changé}
\begin{calculs}
p &= \dfrac{72}{144} &&\justif{$\frac{n}{N}$}\\
p &= 0{,}5\\
p &= 50\,\%
\end{calculs}
\explique{\res{$50\,\%$} des femmes sont mineures, contre $35\,\%$ dans tout le club :
\res{les femmes du club sont plus souvent mineures}.}{on compare deux proportions}
\cas{d)}
\explique{Population de référence : les licenciés du second club, en nombre $N$ inconnu ;
sous-population : les $54$ mineurs.}{c'est $N$ qu'on cherche}
\begin{calculs}
p\times N &= n &&\justif{propriété 2 du cours}\\
0{,}3\times N &= 54\\
N &= \dfrac{54}{0{,}3} &&\justif{on divise les deux membres par $0{,}3$}\\
N &= 180
\end{calculs}
\explique{Ce club compte \res{$180$ licenciés}.}{vérification : $0{,}3\times 180=54$}
\end{exo}

\begin{exo}[Population de référence]
\cas{a)}
\explique{\res{La population de référence est l'ensemble des $2{,}4$ millions de voitures
vendues dans l'année. La sous-population est celle des $360\,000$ voitures
électriques.}}{la sous-population est une partie de la population de référence}
\cas{b)}
\explique{$2{,}4$ millions $=2\,400\,000$.}{on écrit les deux effectifs dans la même unité}
\begin{calculs}
p &= \dfrac{360\,000}{2\,400\,000} &&\justif{$\frac{n}{N}$}\\
p &= 0{,}15\\
p &= 15\,\%
\end{calculs}
\explique{Les voitures électriques représentent \res{$15\,\%$} des ventes.}{on conclut par
une phrase}
\end{exo}

\begin{exo}[Retrouver l'effectif total]
\explique{Dans les trois cas, on connaît la proportion $p$ et l'effectif $n$ de la
sous-population : on cherche $N$, avec $p\times N=n$.}{propriété 2 du cours}
\cas{a)}
\begin{calculs}
0{,}015\times N &= 300 &&\justif{$1{,}5\,\%=0{,}015$}\\
N &= \dfrac{300}{0{,}015} &&\justif{on divise par $0{,}015$}\\
N &= 20\,000
\end{calculs}
\explique{La population compte \res{$20\,000$ insectes}.}{vérification :
$0{,}015\times 20\,000=300$}
\cas{b)}
\begin{calculs}
0{,}6\times N &= 4{,}8 &&\justif{en milliards}\\
N &= \dfrac{4{,}8}{0{,}6}\\
N &= 8
\end{calculs}
\explique{Cela suppose une population mondiale de \res{$8$ milliards} d'habitants.}{on
garde l'unité : des milliards}
\cas{c)}
\begin{calculs}
0{,}6\times N &= 21\\
N &= \dfrac{21}{0{,}6}\\
N &= 35
\end{calculs}
\explique{La classe compte \res{$35$ élèves}.}{vérification : $0{,}6\times 35=21$}
\rem{Erreur fréquente : calculer $60\,\%$ de $21$. Le nombre $21$ est une partie de la
classe, pas la classe entière.}
\end{exo}

\begin{exo}[Un budget]
\cas{a)}
\explique{Le loyer est une partie du salaire $N$, inconnu.}{$p=0{,}4$ et $n=640$}
\begin{calculs}
0{,}4\times N &= 640\\
N &= \dfrac{640}{0{,}4} &&\justif{on divise par $0{,}4$}\\
N &= 1\,600
\end{calculs}
\explique{Le salaire de Yanis est de \res{$1\,600$~€}.}{on conclut par une phrase}
\cas{b)}
\begin{calculs}
\text{charges} &= 0{,}08\times 1\,600 &&\justif{$8\,\%$ du salaire trouvé en a)}\\
\text{charges} &= 128
\end{calculs}
\explique{Les charges s'élèvent à \res{$128$~€}.}{on conclut par une phrase}
\cas{c)}
\begin{calculs}
\dfrac{640+128}{1\,600} &= \dfrac{768}{1\,600} &&\justif{loyer et charges, divisés par le salaire}\\
\dfrac{640+128}{1\,600} &= 0{,}48\\
\dfrac{640+128}{1\,600} &= 48\,\%
\end{calculs}
\explique{Le loyer et les charges représentent \res{$48\,\%$} du salaire.}{plus simple :
$40\,\%+8\,\%=48\,\%$}
\rem{On peut ajouter ces deux pourcentages parce qu'ils portent sur le même total, le
salaire. Des évolutions successives, elles, ne s'ajoutent pas (partie VI).}
\end{exo}

\begin{exo}[Vérifier une information]
\cas{a)}
\begin{calculs}
0{,}03\times N &= 1{,}2 &&\justif{en millions ; $3\,\%$ de la population}\\
N &= \dfrac{1{,}2}{0{,}03}\\
N &= 40
\end{calculs}
\explique{Le pays compterait \res{$40$ millions d'habitants}.}{}
\cas{b)}
\begin{calculs}
\dfrac{1{,}2}{68} &\approx 0{,}0176 &&\justif{calculatrice : $0{,}017647\ldots$}\\
\dfrac{1{,}2}{68} &\approx 1{,}8\,\% &&\justif{arrondi au dixième}
\end{calculs}
\explique{Environ \res{$1{,}8\,\%$} de la population a eu la grippe, et non
$3\,\%$.}{l'affirmation du journal est fausse}
\end{exo}

\partie{Proportions échelonnées}

\begin{exo}[Pourcentage d'un pourcentage]
\explique{Prendre un pourcentage d'un pourcentage, c'est multiplier les deux
proportions.}{propriété 3 du cours}
\begin{calculs}
20\,\%\text{ de }40\,\% &= 0{,}2\times 0{,}4 = 0{,}08\\
20\,\%\text{ de }40\,\% &= \res{8\,\%}\\
50\,\%\text{ de }30\,\% &= 0{,}5\times 0{,}3 = 0{,}15\\
50\,\%\text{ de }30\,\% &= \res{15\,\%}\\
10\,\%\text{ de }5\,\% &= 0{,}1\times 0{,}05 = 0{,}005\\
10\,\%\text{ de }5\,\% &= \res{0{,}5\,\%}
\end{calculs}
\rem{Le résultat est plus petit que chacun des deux pourcentages : on prend une partie
d'une partie.}
\end{exo}

\begin{exo}[Association sportive]
\cas{a)}
\explique{Les joueurs de badminton sont une partie des inscrits, qui sont une partie des
élèves.}{on repère les populations emboîtées}
\begin{calculs}
p &= p_1\times p_2 &&\justif{propriété 3}\\
p &= 0{,}2\times 0{,}45 &&\justif{$20\,\%$ des inscrits, $45\,\%$ des élèves}\\
p &= 0{,}09\\
p &= 9\,\%
\end{calculs}
\explique{\res{$9\,\%$} des élèves du lycée font du badminton à l'association
sportive.}{on conclut par une phrase}
\cas{b)}
\begin{calculs}
n &= p\times N &&\justif{proportion trouvée en a)}\\
n &= 0{,}09\times 1\,200\\
n &= 108
\end{calculs}
\explique{Cela représente \res{$108$ élèves}.}{on conclut par une phrase}
\cas{c)}
\begin{calculs}
\text{inscrits} &= 0{,}45\times 1\,200 &&\justif{$45\,\%$ des élèves}\\
\text{inscrits} &= 540\\
\text{badminton} &= 0{,}2\times 540 &&\justif{$20\,\%$ des inscrits}\\
\text{badminton} &= 108
\end{calculs}
\explique{On retrouve \res{$108$ élèves}.}{les deux méthodes concordent}
\end{exo}

\begin{exo}[Populations emboîtées]
\cas{a)}
\begin{calculs}
p &= 0{,}35\times 0{,}8 &&\justif{pommes $\subset$ fruits $\subset$ compote}\\
p &= 0{,}28\\
p &= 28\,\%
\end{calculs}
\explique{La compote contient \res{$28\,\%$} de pommes.}{on conclut par une phrase}
\cas{b)}
\begin{calculs}
p &= 0{,}4\times 0{,}25 &&\justif{acidulés à la fraise $\subset$ fraise $\subset$ sachet}\\
p &= 0{,}1\\
n &= 0{,}1\times 300 &&\justif{$n=p\times N$}\\
n &= 30
\end{calculs}
\explique{Le sachet contient \res{$30$ bonbons acidulés à la fraise}.}{vérification :
$75$ à la fraise, dont $0{,}4\times 75=30$}
\cas{c)}
\explique{Trois populations emboîtées dans celle des habitants : on multiplie les trois
proportions.}{la propriété se prolonge}
\begin{calculs}
p &= 0{,}25\times 0{,}6\times 0{,}75 &&\justif{du plus petit groupe au plus grand}\\
p &= 0{,}1125\\
p &= 11{,}25\,\%\\
n &= 0{,}1125\times 48\,000 &&\justif{$n=p\times N$}\\
n &= 5\,400
\end{calculs}
\explique{\res{$11{,}25\,\%$} des habitants roulent en voiture électrique, soit
\res{$5\,400$ habitants}.}{une phrase par question}
\end{exo}

\begin{exo}[Tableau croisé]
\cas{a)}
\explique{On calcule d'abord le nombre de filles, puis celui des garçons.}{on part du
total, $600$}
\begin{calculs}
\text{filles} &= 0{,}55\times 600 = 330\\
\text{garçons} &= 600-330 = 270\\
\text{filles droitières} &= 0{,}9\times 330 = 297 &&\justif{$90\,\%$ des filles}\\
\text{filles gauchères} &= 330-297 = 33\\
\text{garçons gauchers} &= 0{,}1\times 270 = 27 &&\justif{$10\,\%$ des garçons}\\
\text{garçons droitiers} &= 270-27 = 243
\end{calculs}
\begin{center}
\renewcommand{\arraystretch}{1.3}
\begin{tabular}{|l|c|c|c|}
\hline
 & Gauchers & Droitiers & Total\\
\hline
Filles & \res{$33$} & \res{$297$} & \res{$330$}\\
\hline
Garçons & \res{$27$} & \res{$243$} & \res{$270$}\\
\hline
Total & \res{$60$} & \res{$540$} & $600$\\
\hline
\end{tabular}
\end{center}
\explique{Vérification : $60+540=600$.}{la dernière ligne s'obtient en additionnant les
colonnes}
\cas{b)}
\begin{calculs}
p &= \dfrac{33}{600} &&\justif{lu dans le tableau}\\
p &= 0{,}055\\
p &= 5{,}5\,\%
\end{calculs}
\explique{Par une proportion échelonnée : les gauchères sont $10\,\%$ des filles, qui sont
$55\,\%$ des élèves.}{$100\,\%-90\,\%=10\,\%$ des filles sont gauchères}
\begin{calculs}
p &= 0{,}1\times 0{,}55\\
p &= 0{,}055
\end{calculs}
\explique{Les filles gauchères représentent \res{$5{,}5\,\%$} des élèves du collège.}{les
deux calculs concordent}
\cas{c)}
\begin{calculs}
p &= \dfrac{60}{600}\\
p &= 0{,}1\\
p &= 10\,\%
\end{calculs}
\explique{\res{$10\,\%$} des élèves du collège sont gauchers.}{on conclut par une phrase}
\end{exo}

\partie{Coefficient multiplicateur}

\begin{exo}[Taux et coefficient multiplicateur]
\cas{a)}
\explique{$C_m=1+t$, où le taux $t$ est positif pour une hausse ($t=0{,}15$ pour
$+15\,\%$), négatif pour une baisse ($t=-0{,}04$ pour $-4\,\%$).}{propriété 4 du cours}
\begin{calculs}
C_m &= 1+0{,}15 = \res{1{,}15} &&\justif{hausse de $15\,\%$}\\
C_m &= 1-0{,}04 = \res{0{,}96} &&\justif{baisse de $4\,\%$}\\
C_m &= 1+1{,}2 = \res{2{,}2} &&\justif{hausse de $120\,\%$ : $\frac{120}{100}=1{,}2$}\\
C_m &= 1-0{,}6 = \res{0{,}4} &&\justif{baisse de $60\,\%$}\\
C_m &= 1-0{,}123 = \res{0{,}877} &&\justif{baisse de $12{,}3\,\%$ : $\frac{12{,}3}{100}=0{,}123$}
\end{calculs}
\cas{b)}
\explique{On écrit $C_m$ sous la forme $1+t$ : le taux $t$ est positif (hausse) ou
négatif (baisse).}{méthode 5 du cours}
\begin{calculs}
C_m &= 1{,}03 = 1+0{,}03 &&\justif{taux d'évolution : \res{$+3\,\%$}}\\
C_m &= 0{,}72 = 1-0{,}28 &&\justif{taux d'évolution : \res{$-28\,\%$}}\\
C_m &= 2 = 1+1 &&\justif{taux d'évolution : \res{$+100\,\%$} (doubler)}\\
C_m &= 0{,}995 = 1-0{,}005 &&\justif{taux d'évolution : \res{$-0{,}5\,\%$}}
\end{calculs}
\rem{Erreur fréquente : lire $0{,}72$ comme une baisse de $72\,\%$. C'est ce qui
\emph{reste} : $72\,\%$ de la valeur, donc une baisse de $28\,\%$.}
\end{exo}

\begin{exo}[En phrases]
\begin{calculs}
1{,}08 &= 1+0{,}08 &&\justif{\res{Multiplier par $1{,}08$, c'est augmenter de $8\,\%$.}}\\
0{,}92 &= 1-0{,}08 &&\justif{\res{Multiplier par $0{,}92$, c'est diminuer de $8\,\%$.}}\\
1{,}5 &= 1+0{,}5 &&\justif{\res{Multiplier par $1{,}5$, c'est augmenter de $50\,\%$.}}\\
0{,}3 &= 1-0{,}7 &&\justif{\res{Multiplier par $0{,}3$, c'est diminuer de $70\,\%$.}}
\end{calculs}
\rem{Un coefficient plus grand que $1$ traduit une hausse ; plus petit que $1$, une
baisse.}
\end{exo}

\begin{exo}[Valeur d'arrivée]
\explique{Valeur d'arrivée : $V_1=V_0\times C_m$.}{propriété 4 du cours}
\cas{a)}
\begin{calculs}
V_1 &= V_0\times C_m\\
V_1 &= 24\,000\times 1{,}05 &&\justif{hausse de $5\,\%$ : $C_m=1+0{,}05$}\\
V_1 &= 25\,200
\end{calculs}
\explique{La ville compte \res{$25\,200$ habitants}.}{on conclut par une phrase}
\cas{b)}
\begin{calculs}
V_1 &= 75\times 0{,}8 &&\justif{réduction de $20\,\%$ : $C_m=1-0{,}2$}\\
V_1 &= 60
\end{calculs}
\explique{Les chaussures coûtent \res{$60$~€}.}{on conclut par une phrase}
\cas{c)}
\begin{calculs}
V_1 &= 3\,500\times 0{,}88 &&\justif{baisse de $12\,\%$ : $C_m=1-0{,}12$}\\
V_1 &= 3\,080
\end{calculs}
\explique{\res{$3\,080$ accidents} ont eu lieu l'année suivante.}{on conclut par une phrase}
\end{exo}

\begin{exo}[Valeur de départ]
\explique{On cherche $V_0$ : on écrit $V_0\times C_m=V_1$, puis on divise par
$C_m$.}{méthode 6 du cours}
\cas{a)}
\begin{calculs}
V_0\times C_m &= V_1\\
V_0\times 1{,}08 &= 27 &&\justif{hausse de $8\,\%$ : $C_m=1+0{,}08$}\\
V_0 &= \dfrac{27}{1{,}08} &&\justif{on divise par $C_m$}\\
V_0 &= 25
\end{calculs}
\explique{L'abonnement coûtait \res{$25$~€}.}{vérification : $25\times 1{,}08=27$}
\cas{b)}
\begin{calculs}
V_0\times 0{,}95 &= 19 &&\justif{baisse de $5\,\%$ : $C_m=1-0{,}05$}\\
V_0 &= \dfrac{19}{0{,}95}\\
V_0 &= 20
\end{calculs}
\explique{La température moyenne d'août était de \res{$20$~°C}.}{vérification :
$20\times 0{,}95=19$}
\cas{c)}
\begin{calculs}
V_0\times 1{,}2 &= 150 &&\justif{la TVA est une hausse de $20\,\%$ du prix HT}\\
V_0 &= \dfrac{150}{1{,}2}\\
V_0 &= 125
\end{calculs}
\explique{Le lit coûte \res{$125$~€} HT.}{retirer $20\,\%$ de $150$ donnerait $120$ : faux}
\cas{d)}
\explique{Ici, on connaît l'augmentation ($27$~€), pas le nouveau prix : $27$~€ sont
$18\,\%$ du prix de départ.}{c'est la partie II : $p\times N=n$}
\begin{calculs}
0{,}18\times V_0 &= 27\\
V_0 &= \dfrac{27}{0{,}18}\\
V_0 &= 150
\end{calculs}
\explique{Le vélo coûtait \res{$150$~€} avant la hausse.}{il coûte maintenant
$150+27=177$~€, et $150\times 1{,}18=177$}
\end{exo}

\begin{exo}[Les soldes]
\cas{a)}
\begin{calculs}
V_1 &= 225\times 0{,}8 &&\justif{réduction de $20\,\%$ : $C_m=1-0{,}2$}\\
V_1 &= 180
\end{calculs}
\explique{Léo aurait payé \res{$180$~€}.}{on cherche la valeur d'arrivée}
\cas{b)}
\begin{calculs}
V_0\times 0{,}8 &= 240 &&\justif{on cherche la valeur de départ}\\
V_0 &= \dfrac{240}{0{,}8}\\
V_0 &= 300
\end{calculs}
\explique{Inès aurait payé \res{$300$~€}.}{ajouter $20\,\%$ à $240$ donnerait $288$ : faux}
\end{exo}

\partie{Calculer un taux d'évolution}

\begin{exo}[Taux d'évolution]
\explique{Variation absolue : $V_1-V_0$ ; taux d'évolution :
$t=\dfrac{V_1-V_0}{V_0}$.}{propriété 5 du cours}
\cas{a)}
\begin{calculs}
V_1-V_0 &= 3\,040-3\,200 &&\justif{$V_0=3\,200$ (départ) et $V_1=3\,040$ (arrivée)}\\
V_1-V_0 &= -160\\
t &= \dfrac{-160}{3\,200} &&\justif{on divise par la valeur de départ}\\
t &= -0{,}05\\
t &= -5\,\%
\end{calculs}
\explique{La commune perd \res{$160$ habitants} : sa population baisse de
\res{$5\,\%$}.}{$t$ négatif : une baisse}
\cas{b)}
\begin{calculs}
V_1-V_0 &= 2\,100-1\,500\\
V_1-V_0 &= 600\\
t &= \dfrac{600}{1\,500}\\
t &= 0{,}4\\
t &= 40\,\%
\end{calculs}
\explique{La production augmente de \res{$600$ vélos}, soit \res{$40\,\%$}.}{$t$ positif :
une hausse}
\cas{c)}
\begin{calculs}
V_1-V_0 &= 1\,782-1\,650\\
V_1-V_0 &= 132\\
t &= \dfrac{132}{1\,650}\\
t &= 0{,}08\\
t &= 8\,\%
\end{calculs}
\explique{Le salaire augmente de \res{$132$~€}, soit \res{$8\,\%$}.}{on conclut par une
phrase}
\cas{d)}
\begin{calculs}
V_1-V_0 &= 19{,}1-18{,}4\\
V_1-V_0 &= 0{,}7\\
t &= \dfrac{0{,}7}{18{,}4}\\
t &\approx 0{,}0380 &&\justif{calculatrice : $0{,}038043\ldots$}\\
t &\approx 3{,}8\,\% &&\justif{arrondi au dixième de pourcentage}
\end{calculs}
\explique{Le plein augmente de \res{$0{,}70$~€}, soit environ \res{$3{,}8\,\%$}.}{on
conclut par une phrase}
\end{exo}

\begin{exo}[Deux façons de dire]
\cas{a)}
\begin{calculs}
V_1-V_0 &= 460-400 &&\justif{variation absolue}\\
V_1-V_0 &= 60\\
t &= \dfrac{60}{400} &&\justif{taux d'évolution}\\
t &= 0{,}15\\
t &= 15\,\%
\end{calculs}
\cas{b)}
\explique{\res{Le club a gagné $60$ adhérents.}}{la variation absolue a l'unité de la
grandeur : des adhérents}
\explique{\res{Le nombre d'adhérents du club a augmenté de $15\,\%$.}}{le taux n'a pas
d'unité : un pourcentage}
\end{exo}

\begin{exo}[Du coefficient au taux]
\cas{a)}
\begin{calculs}
C_m &= \dfrac{V_1}{V_0} &&\justif{rappel de l'énoncé}\\
C_m &= \dfrac{144}{120}\\
C_m &= 1{,}2\\
t &= C_m-1\\
t &= 1{,}2-1\\
t &= 0{,}2
\end{calculs}
\explique{$C_m=\res{1{,}2}$ et $t=\res{20\,\%}$ : le prix augmente de $20\,\%$.}{même
résultat que $\frac{144-120}{120}$}
\cas{b)}
\begin{calculs}
C_m &= \dfrac{200}{250}\\
C_m &= 0{,}8\\
t &= 0{,}8-1\\
t &= -0{,}2
\end{calculs}
\explique{$C_m=\res{0{,}8}$ et $t=\res{-20\,\%}$ : le stock baisse de $20\,\%$.}{$t$
négatif : une baisse}
\cas{c)}
\begin{calculs}
t &= C_m-1\\
t &= 1{,}5-1\\
t &= 0{,}5
\end{calculs}
\explique{$t=\res{50\,\%}$ : le prix augmente de $50\,\%$.}{}
\cas{d)}
\begin{calculs}
t &= 0{,}65-1\\
t &= -0{,}35
\end{calculs}
\explique{$t=\res{-35\,\%}$ : le prix baisse de $35\,\%$.}{même lecture que $0{,}65=1-0{,}35$}
\end{exo}

\begin{exo}[Vrai ou faux]
\cas{a)}
\begin{calculs}
t &= \dfrac{150-120}{120} &&\justif{on divise par la valeur de départ, $120$}\\
t &= \dfrac{30}{120}\\
t &= 0{,}25
\end{calculs}
\explique{\res{Faux} : le prix augmente de $25\,\%$. L'augmentation vaut $30$~€, pas
$30\,\%$.}{ne pas confondre variation absolue et taux}
\cas{b)}
\begin{calculs}
C_m &= \dfrac{150}{120}\\
C_m &= 1{,}25
\end{calculs}
\explique{\res{Vrai}.}{cohérent avec a) : $1{,}25=1+0{,}25$}
\cas{c)}
\begin{calculs}
t &= \dfrac{200-250}{250}\\
t &= \dfrac{-50}{250}\\
t &= -0{,}2
\end{calculs}
\explique{\res{Faux} : le stock baisse de $20\,\%$.}{diviser par $200$ (la valeur
d'arrivée) donnerait $25\,\%$ : c'est l'erreur}
\cas{d)}
\begin{calculs}
t &= \dfrac{-60}{300}\\
t &= -0{,}2
\end{calculs}
\explique{\res{Vrai} : la population baisse de $20\,\%$.}{la variation absolue est $-60$}
\end{exo}

\begin{exo}[Deux magasins]
\cas{a)}
\begin{calculs}
t &= \dfrac{408-480}{480} &&\justif{magasin A}\\
t &= \dfrac{-72}{480}\\
t &= -0{,}15
\end{calculs}
\explique{Le magasin A baisse son prix de $15\,\%$, le magasin B de $14\,\%$ : la plus forte
baisse est dans le \res{magasin A}.}{on compare les taux}
\cas{b)}
\begin{calculs}
V_1 &= 500\times 0{,}86 &&\justif{magasin B : baisse de $14\,\%$, $C_m=1-0{,}14$}\\
V_1 &= 430
\end{calculs}
\explique{Le téléviseur coûte $408$~€ dans le magasin A et $430$~€ dans le B : il est le
moins cher dans le \res{magasin A}.}{on compare les prix}
\end{exo}

\begin{exo}[Une entreprise qui embauche]
\cas{a)}
\begin{calculs}
\text{salariés} &= 250+30\\
\text{salariés} &= 280
\end{calculs}
\explique{L'entreprise compte \res{$280$ salariés}.}{}
\cas{b)}
\begin{calculs}
\dfrac{30}{250} &= 0{,}12 &&\justif{les embauches, divisées par l'effectif de départ}\\
\dfrac{30}{250} &= 12\,\%
\end{calculs}
\explique{Les embauches représentent \res{$12\,\%$} de l'effectif de départ.}{c'est aussi
le taux d'évolution de la première année}
\cas{c)}
\begin{calculs}
\text{salariés} &= 280+14\\
\text{salariés} &= 294
\end{calculs}
\explique{L'entreprise compte \res{$294$ salariés}.}{}
\cas{d)}
\begin{calculs}
t &= \dfrac{294-250}{250} &&\justif{de $V_0=250$ à $V_1=294$}\\
t &= \dfrac{44}{250}\\
t &= 0{,}176\\
t &= 17{,}6\,\%
\end{calculs}
\explique{Le nombre de salariés augmente de \res{$17{,}6\,\%$} en deux ans.}{on divise
toujours par la valeur de départ, $250$}
\end{exo}

\partie{Évolutions successives}

\begin{exo}[Taux global]
\explique{Le coefficient multiplicateur global est le produit des coefficients :
$C_m=C_{m1}\times C_{m2}$.}{propriété 6 du cours}
\cas{a)}
\begin{calculs}
C_m &= 1{,}25\times 1{,}1 &&\justif{$1+0{,}25$, puis $1+0{,}1$}\\
C_m &= 1{,}375\\
C_m &= 1+0{,}375
\end{calculs}
\explique{Hausse globale de \res{$37{,}5\,\%$} (et non de $35\,\%$).}{$0{,}375=37{,}5\,\%$}
\cas{b)}
\begin{calculs}
C_m &= 1{,}3\times 0{,}7 &&\justif{$1+0{,}3$, puis $1-0{,}3$}\\
C_m &= 0{,}91\\
C_m &= 1-0{,}09
\end{calculs}
\explique{Baisse globale de \res{$9\,\%$} : les deux évolutions ne se compensent
pas.}{$0{,}09=9\,\%$}
\cas{c)}
\begin{calculs}
C_m &= 0{,}6\times 0{,}9 &&\justif{$1-0{,}4$, puis $1-0{,}1$}\\
C_m &= 0{,}54\\
C_m &= 1-0{,}46
\end{calculs}
\explique{Baisse globale de \res{$46\,\%$} (et non de $50\,\%$).}{$0{,}46=46\,\%$}
\cas{d)}
\begin{calculs}
C_m &= 1{,}02\times 0{,}99 &&\justif{$1+0{,}02$, puis $1-0{,}01$}\\
C_m &= 1{,}0098\\
C_m &= 1+0{,}0098
\end{calculs}
\explique{Hausse globale de \res{$0{,}98\,\%$}.}{$0{,}0098=0{,}98\,\%$}
\end{exo}

\begin{exo}[Un ticket de transport]
\cas{a)}
\begin{calculs}
C_m &= C_{m1}\times C_{m2} &&\justif{propriété 6}\\
C_m &= 1{,}3\times 0{,}9 &&\justif{hausse de $30\,\%$, puis baisse de $10\,\%$}\\
C_m &= 1{,}17\\
C_m &= 1+0{,}17
\end{calculs}
\explique{Le taux d'évolution global est de \res{$+17\,\%$} (et non de $+20\,\%$).}{}
\cas{b)}
\begin{calculs}
V_2 &= V_0\times C_m &&\justif{avec le coefficient global}\\
V_2 &= 25\times 1{,}17\\
V_2 &= 29{,}25
\end{calculs}
\explique{L'abonnement coûte \res{$29{,}25$~€}.}{vérification : $25\times 1{,}3=32{,}5$, puis
$32{,}5\times 0{,}9=29{,}25$}
\end{exo}

\begin{exo}[Vrai ou faux]
\cas{a)}
\begin{calculs}
C_m &= 1{,}12\times 0{,}88\\
C_m &= 0{,}9856\\
C_m &= 1-0{,}0144
\end{calculs}
\explique{\res{Faux} : le prix baisse d'environ $1{,}4\,\%$.}{la baisse porte sur un prix
déjà augmenté}
\cas{b)}
\begin{calculs}
C_m &= 0{,}93\times 0{,}8 = 0{,}744 &&\justif{$7\,\%$ puis $20\,\%$}\\
C_m &= 0{,}8\times 0{,}93 = 0{,}744 &&\justif{$20\,\%$ puis $7\,\%$}
\end{calculs}
\explique{\res{Vrai} : l'ordre d'une multiplication ne change pas le résultat.}{les deux
coefficients sont les mêmes}
\cas{c)}
\begin{calculs}
C_m &= 1{,}06\times 1{,}08\\
C_m &= 1{,}1448\\
C_m &= 1+0{,}1448
\end{calculs}
\explique{\res{Faux} : la hausse est de $14{,}48\,\%$, un peu plus que $14\,\%$.}{la seconde
hausse porte sur un prix déjà augmenté}
\cas{d)}
\begin{calculs}
C_m &= 1{,}1\times 1{,}15\\
C_m &= 1{,}265\\
C_m &= 1+0{,}265
\end{calculs}
\explique{\res{Faux} : le prix augmente de $26{,}5\,\%$, plus qu'un quart ($25\,\%$).}{on
n'ajoute pas les pourcentages}
\end{exo}

\begin{exo}[Le fleuriste]
\begin{calculs}
C_m &= 1{,}4\times 0{,}7 &&\justif{hausse de $40\,\%$, puis réduction de $30\,\%$}\\
C_m &= 0{,}98\\
V_2 &= 8\times 0{,}98 &&\justif{prix de vente final d'un pot}\\
V_2 &= 7{,}84
\end{calculs}
\explique{Le fleuriste vend $7{,}84$~€ un pot acheté $8$~€ : \res{oui, il vend à
perte}.}{$C_m=0{,}98$, plus petit que $1$ : une baisse globale de $2\,\%$}
\end{exo}

\begin{exo}[La valeur d'une voiture]
\cas{a)}
\begin{calculs}
V_1 &= 18\,000\times 0{,}8 &&\justif{baisse de $20\,\%$ : $C_m=1-0{,}2$}\\
V_1 &= 14\,400\\
V_2 &= 14\,400\times 0{,}85 &&\justif{baisse de $15\,\%$ : $C_m=1-0{,}15$}\\
V_2 &= 12\,240
\end{calculs}
\explique{La voiture vaut \res{$14\,400$~€} au bout d'un an, \res{$12\,240$~€} au bout de
deux ans.}{}
\cas{b)}
\begin{calculs}
V_3 &= 12\,240\times 0{,}85\\
V_3 &= 10\,404
\end{calculs}
\explique{Elle vaut \res{$10\,404$~€} au bout de trois ans.}{}
\cas{c)}
\begin{calculs}
C_m &= 0{,}8\times 0{,}85\times 0{,}85 &&\justif{trois évolutions successives}\\
C_m &= 0{,}578\\
C_m &= 1-0{,}422
\end{calculs}
\explique{Le taux d'évolution global est d'environ \res{$-42{,}2\,\%$}.}{vérification :
$18\,000\times 0{,}578=10\,404$}
\end{exo}

\partie{Évolution réciproque}

\begin{exo}[Revenir au départ]
\explique{Le coefficient de l'évolution réciproque est l'inverse :
$C'_m=\dfrac{1}{C_m}$.}{propriété 7 du cours}
\cas{a)}
\begin{calculs}
C'_m &= \dfrac{1}{0{,}8} &&\justif{baisse de $20\,\%$ : $C_m=1-0{,}2$}\\
C'_m &= 1{,}25\\
C'_m &= 1+0{,}25
\end{calculs}
\explique{Il faut une \res{hausse de $25\,\%$}.}{et non de $20\,\%$}
\cas{b)}
\begin{calculs}
C'_m &= \dfrac{1}{2} &&\justif{hausse de $100\,\%$ : $C_m=1+1$}\\
C'_m &= 0{,}5\\
C'_m &= 1-0{,}5
\end{calculs}
\explique{Il faut une \res{baisse de $50\,\%$} : diviser par deux compense le
doublement.}{}
\cas{c)}
\begin{calculs}
C'_m &= \dfrac{1}{0{,}55} &&\justif{baisse de $45\,\%$ : $C_m=1-0{,}45$}\\
C'_m &\approx 1{,}8182 &&\justif{calculatrice : $1{,}818181\ldots$}\\
C'_m &\approx 1+0{,}8182
\end{calculs}
\explique{Il faut une \res{hausse d'environ $81{,}8\,\%$}.}{}
\cas{d)}
\begin{calculs}
C'_m &= \dfrac{1}{1{,}3} &&\justif{hausse de $30\,\%$ : $C_m=1+0{,}3$}\\
C'_m &\approx 0{,}7692 &&\justif{calculatrice : $0{,}769230\ldots$}\\
C'_m &\approx 1-0{,}2308
\end{calculs}
\explique{Il faut une \res{baisse d'environ $23{,}1\,\%$}.}{}
\end{exo}

\begin{exo}[Un village]
\cas{a)}
\begin{calculs}
t &= \dfrac{3\,000-2\,500}{2\,500}\\
t &= \dfrac{500}{2\,500}\\
t &= 0{,}2
\end{calculs}
\explique{La population augmente de \res{$20\,\%$}.}{}
\cas{b)}
\begin{calculs}
C'_m &= \dfrac{1}{1{,}2} &&\justif{hausse de $20\,\%$ : $C_m=1+0{,}2$}\\
C'_m &\approx 0{,}8333 &&\justif{calculatrice : $0{,}833333\ldots$}\\
C'_m &\approx 1-0{,}1667
\end{calculs}
\explique{Une \res{baisse d'environ $16{,}7\,\%$} ramènerait le village à $2\,500$
habitants.}{autre voie : $\frac{2\,500-3\,000}{3\,000}\approx -0{,}1667$}
\end{exo}

\begin{exo}[Aller et retour]
\cas{a)}
\begin{calculs}
t &= \dfrac{26-20}{20} = 0{,}3 &&\justif{taux d'évolution : $+30\,\%$}\\
C_m &= 1+0{,}3 = 1{,}3\\
C'_m &= \dfrac{1}{1{,}3} \approx 0{,}7692 &&\justif{inverse de $C_m$}\\
C'_m &\approx 1-0{,}2308
\end{calculs}
\explique{Taux : \res{$+30\,\%$} ; $C_m=\res{1{,}3}$ ; évolution réciproque : \res{baisse
d'environ $23{,}1\,\%$}.}{}
\cas{b)}
\begin{calculs}
t &= \dfrac{39-52}{52} = -0{,}25 &&\justif{taux d'évolution : $-25\,\%$}\\
C_m &= 1-0{,}25 = 0{,}75\\
C'_m &= \dfrac{1}{0{,}75} \approx 1{,}3333\\
C'_m &\approx 1+0{,}3333
\end{calculs}
\explique{Taux : \res{$-25\,\%$} ; $C_m=\res{0{,}75}$ ; évolution réciproque : \res{hausse
d'environ $33{,}3\,\%$}.}{vérification : $39\times 1{,}3333\approx 52$}
\end{exo}

\begin{exo}[Plus de chocolat pour le même prix]
\explique{La quantité est multipliée par $1{,}25$ pour le même prix : le prix du kilo est
divisé par $1{,}25$, c'est-à-dire multiplié par $\dfrac{1}{1{,}25}$.}{c'est l'évolution
réciproque d'une hausse de $25\,\%$}
\begin{calculs}
C'_m &= \dfrac{1}{1{,}25}\\
C'_m &= 0{,}8\\
C'_m &= 1-0{,}2
\end{calculs}
\explique{Le prix du kilo de chocolat a baissé de \res{$20\,\%$} (et non de
$25\,\%$).}{exemple : $10$~€ pour $1$~kg, puis $10$~€ pour $1{,}25$~kg, soit $8$~€ le kilo}
\end{exo}

\begin{exo}[Retrouver un prix]
\cas{a)}
\begin{calculs}
C_m &= 1{,}1\times 0{,}8 &&\justif{hausse de $10\,\%$, puis baisse de $20\,\%$}\\
C_m &= 0{,}88\\
C_m &= 1-0{,}12
\end{calculs}
\explique{$C_m=\res{0{,}88}$ : baisse globale de \res{$12\,\%$}.}{}
\cas{b)}
\begin{calculs}
V_0\times C_m &= V_2 &&\justif{on cherche la valeur de départ}\\
V_0\times 0{,}88 &= 352\\
V_0 &= \dfrac{352}{0{,}88}\\
V_0 &= 400
\end{calculs}
\explique{Le casque coûtait \res{$400$~€}.}{vérification : $400\times 1{,}1=440$, puis
$440\times 0{,}8=352$}
\cas{c)}
\begin{calculs}
C'_m &= \dfrac{1}{0{,}88}\\
C'_m &\approx 1{,}1364 &&\justif{calculatrice : $1{,}136363\ldots$}\\
C'_m &\approx 1+0{,}1364
\end{calculs}
\explique{Une \res{hausse d'environ $13{,}6\,\%$} ramènerait le prix à $400$~€.}{et non de
$12\,\%$}
\end{exo}

\end{document}
